\documentclass[11pt]{article}
\usepackage[margin=1in]{geometry}
\usepackage{amsmath, amssymb}
\usepackage{booktabs}
\usepackage{graphicx}
\usepackage{hyperref}
\usepackage{caption}
\usepackage{subcaption}
\usepackage{multirow}
\usepackage{array}

\title{Detecting Adversarial Perturbations in AI-Generated Code: A Study of Static and Embedding-Based Approaches}

\author{Anonymous Author(s)}

\date{}

\begin{document}
\maketitle

\begin{abstract}
AI code generation systems are increasingly used to produce production software, yet the security implications of adversarially perturbed AI-generated code remain underexplored. We systematically study five perturbation strategies that an adversary could use to introduce subtle bugs or vulnerabilities into AI-generated Python functions while preserving syntactic plausibility. Using a dataset of 550 samples (100 clean, 450 perturbed), we evaluate two detection approaches: static AST-based features with gradient-boosted trees, and CodeBERT embeddings with logistic regression. We find that AST-based features achieve 91\% accuracy and 96\% precision but exhibit a fundamental blind spot: they detect \textit{structural} perturbations (comment planting, dead code insertion, variable shadowing) with 100\% recall, yet fail completely on \textit{semantic} perturbations (import aliasing, boundary inversion) that preserve syntactic structure while changing program behavior. CodeBERT embeddings improve with scale but do not close this gap. Our findings reveal a three-tier detectability hierarchy and highlight an urgent need for execution-aware or token-level detection methods to secure AI-generated code.
\end{abstract}

\section{Introduction}
\label{sec:intro}

Large language models (LLMs) are increasingly deployed as coding assistants, with tools like GitHub Copilot, Codex, and Gemini generating production code at scale. While much attention has focused on accidental bugs in AI-generated code, the \textit{adversarial} threat model remains critically underexplored: what if an attacker could systematically induce vulnerabilities in AI-generated code that pass both human review and automated static analysis?

We address this gap by asking three research questions:

\begin{enumerate}
    \item What perturbation strategies can adversarially modify AI-generated code while preserving syntactic plausibility?
    \item How effectively do current detection approaches—static AST features and learned embeddings—identify these perturbations?
    \item Is there a fundamental limit to what static and embedding-based methods can detect?
\end{enumerate}

Our contributions are:
\begin{itemize}
    \item A systematic taxonomy of five perturbation strategies (Section~\ref{sec:method})
    \item A curated dataset of 550 AI-generated Python functions with labeled perturbations (Section~\ref{sec:dataset})
    \item A comparative evaluation of AST-based and embedding-based detectors (Section~\ref{sec:results})
    \item The discovery of a three-tier detectability hierarchy revealing that semantic perturbations are invisible to current static methods (Section~\ref{sec:discussion})
\end{itemize}

\section{Related Work}
\label{sec:related}

Prior work on AI code security has focused primarily on \textit{accidental} vulnerabilities. Pearce et al.~\cite{pearce2022asleep} studied security bugs in Copilot-generated code, finding that approximately 40\% of generated functions contain vulnerabilities. Chen et al.~\cite{chen2021evaluating} proposed Codex-based vulnerability detection. However, these works assume bugs arise from model limitations, not from adversarial intent.

In adversarial machine learning, perturbation attacks on NLP systems are well-studied~\cite{jia2017adversarial}. Yet the transfer of these ideas to code generation remains nascent. Recent work on adversarial code examples~\cite{yang2022natural} explores evasion of code classification models, but the reverse problem—detecting adversarially perturbed code—has received limited attention.

Our work bridges this gap by systematically characterizing the detectability limits of current methods.

\section{Methodology}
\label{sec:method}

\subsection{Threat Model}
We assume an adversary with the ability to modify AI-generated code after generation but before deployment. The adversary's goal is to introduce a vulnerability or logic bug that:
\begin{itemize}
    \item Preserves syntactic plausibility (passes human review)
    \item Evades static analysis tools
    \item Alters program behavior in a harmful way
\end{itemize}

\subsection{Perturbation Strategies}
\label{sec:strategies}

We define five perturbation strategies:

\begin{table}[h]
\centering
\begin{tabular}{@{}lll@{}}
\toprule
Strategy & Type & Description \\
\midrule
Comment planting & Structural & Insert commented-out malicious code \\
Dead code insertion & Structural & Add \texttt{if False:} blocks with payloads \\
Variable shadowing & Structural & Reassign parameters to mask original values \\
Import aliasing & Semantic & Alias unsafe imports with misleading names \\
Boundary inversion & Semantic & Flip comparison operators or boundary conditions \\
\bottomrule
\end{tabular}
\caption{The five perturbation strategies studied. Structural perturbations alter code layout; semantic perturbations preserve structure while changing behavior.}
\label{tab:strategies}
\end{table}

\subsection{Detection Approaches}
\label{sec:detectors}

\textbf{Track A: AST Features + XGBoost.} We extract 17 structural features from each code sample's abstract syntax tree (AST), including node counts, depth, import patterns, and suspicious-pattern flags. We train an XGBoost classifier with scale-positive-weight balancing.

\textbf{Track B: CodeBERT Embeddings + Logistic Regression.} We embed each code sample using CodeBERT~\cite{feng2020codebert}, mean-pooling the final hidden states to produce 768-dimensional vectors. We train a logistic regression classifier with balanced class weights.

\textbf{Ensemble.} We concatenate AST features and embeddings (785 dimensions total) and train XGBoost on the combined representation.

\section{Dataset}
\label{sec:dataset}

We create a dataset of 100 clean AI-generated Python functions covering common programming tasks (string manipulation, list operations, mathematical computations, data structures). For each clean function, we automatically apply all five perturbation strategies, yielding 450 perturbed samples. The total dataset contains 550 samples (100 clean, 450 perturbed).

\section{Results}
\label{sec:results}

\subsection{Overall Detection Performance}

\begin{table}[h]
\centering
\begin{tabular}{@{}lccc@{}}
\toprule
Detector & Accuracy & Precision & Recall \\
\midrule
AST Features + XGBoost & \textbf{0.91} $\pm$ 0.03 & \textbf{0.96} $\pm$ 0.04 & 0.93 $\pm$ 0.02 \\
CodeBERT + LogReg & 0.87 $\pm$ 0.03 & 0.91 $\pm$ 0.02 & \textbf{0.93} $\pm$ 0.03 \\
Ensemble & 0.87 $\pm$ 0.03 & 0.92 $\pm$ 0.02 & 0.92 $\pm$ 0.02 \\
\bottomrule
\end{tabular}
\caption{Cross-validated detection performance (5-fold). AST features outperform embeddings and ensemble approaches.}
\label{tab:overall}
\end{table}

\subsection{Leave-One-Strategy-Out Analysis}
\label{sec:loso}

Our most significant finding emerges from leave-one-strategy-out evaluation, where we train on all strategies except one and evaluate on the held-out strategy.

\begin{table}[h]
\centering
\begin{tabular}{@{}lccc@{}}
\toprule
Held-out Strategy & N & Accuracy & Recall \\
\midrule
Comment planting & 100 & 1.00 & 1.00 \\
Dead code insertion & 100 & 1.00 & 1.00 \\
Variable shadowing & 100 & 1.00 & 1.00 \\
Import aliasing & 100 & 0.07 & 0.07 \\
Boundary inversion & 40 & 0.00 & 0.00 \\
\bottomrule
\end{tabular}
\caption{Leave-one-strategy-out results for AST-based detector. Structural perturbations are perfectly detected; semantic perturbations are invisible.}
\label{tab:loso}
\end{table}

\section{Discussion}
\label{sec:discussion}

\subsection{Three-Tier Detectability Hierarchy}

Our results reveal a clear hierarchy in perturbation detectability:

\begin{itemize}
    \item \textbf{Tier 1: Trivially detectable.} Comment planting, dead code insertion, and variable shadowing alter code structure in ways that AST features capture perfectly (recall = 1.00).
    \item \textbf{Tier 2: Undetectable by AST, partially detectable by embeddings.} Import aliasing is caught when the \texttt{has\_aliased\_import} feature is present, but generalizes poorly (recall = 0.07 in leave-one-out).
    \item \textbf{Tier 3: Invisible.} Boundary inversion flips a single operator, producing no structural trace. Neither AST features nor embeddings detect it (recall = 0.00).
\end{itemize}

\subsection{Why Embeddings Don't Help}

CodeBERT embeddings improve from 0.79 accuracy (N=146) to 0.87 (N=550), confirming that embeddings benefit from scale. However, mean-pooling washes out single-token changes (operator flips, import aliases), which carry the semantic signal. Token-level attention or execution-aware analysis may be necessary.

\subsection{Implications for AI Code Security}

Our findings have urgent implications: semantic perturbations represent a viable attack vector against AI-generated code that passes both human review and static analysis. As AI coding assistants become ubiquitous, the development of execution-aware or token-level detection methods becomes critical.

\section{Limitations}

\begin{itemize}
    \item Dataset uses synthetic perturbations; real-world attackers may use more sophisticated strategies.
    \item Python-only; results may not transfer to other languages.
    \item Mean-pooling for embeddings may not be optimal; future work should explore token-level methods.
    \item No execution-based analysis; static methods cannot observe behavioral differences.
\end{itemize}

\section{Conclusion}

We present the first systematic study of adversarial perturbation detection in AI-generated code. Our three-tier detectability hierarchy reveals that current static and embedding-based methods fail catastrophically on semantic perturbations that preserve syntactic structure. This work establishes a foundation for future research on execution-aware code security and motivates the development of token-level detection systems.

\section*{Reproducibility Statement}

All code is available at \texttt{github.com/anonymous/adversarial-code-perturbations}. Experiments require Python 3.14, PyTorch, Transformers, XGBoost, and scikit-learn. The dataset is generated deterministically from 100 seed functions.

\bibliographystyle{plain}
\begin{thebibliography}{9}

\bibitem{pearce2022asleep}
H. Pearce, B. Ahmad, B. Tan, B. Dolan-Gavitt, and R. Karri.
\newblock Asleep at the keyboard? Assessing the security of GitHub Copilot's code contributions.
\newblock In \textit{IEEE Symposium on Security and Privacy (S\&P)}, 2022.

\bibitem{chen2021evaluating}
M. Chen et al.
\newblock Evaluating large language models trained on code.
\newblock \textit{arXiv preprint arXiv:2107.03374}, 2021.

\bibitem{jia2017adversarial}
R. Jia and P. Liang.
\newblock Adversarial examples for evaluating reading comprehension systems.
\newblock In \textit{EMNLP}, 2017.

\bibitem{yang2022natural}
Z. Yang et al.
\newblock Natural adversarial examples for code.
\newblock In \textit{Findings of ACL}, 2022.

\bibitem{feng2020codebert}
Z. Feng et al.
\newblock CodeBERT: A pre-trained model for programming and natural languages.
\newblock In \textit{Findings of EMNLP}, 2020.

\end{thebibliography}

\end{document}
